\documentclass[a4paper, serif, 10pt]{moderncv}

%% moderncv
% color and style
\moderncvcolor{black}
\moderncvstyle{banking}

%% page layout/margins
\usepackage[top=2.5cm, bottom=2.5cm, left=1.5cm, right=1.5cm]{geometry}

\nopagenumbers{}

%% Babel config for English language
\usepackage[english]{babel}

%% fontspec
\usepackage{fontspec}

\IfFontExistsTF{Linux Libertine}{
  \setmainfont[Ligatures={TeX, Common}, Numbers=OldStyle]{Linux Libertine}
}{}
\IfFontExistsTF{Linux Libertine O}{
  \setmainfont[Ligatures={TeX, Common}, Numbers=OldStyle]{Linux Libertine O}
}{}

%% CTeX
% CJK support, used to show CJK characters in the resume
%
% - heading=false: disable ctex heading
% - punct=kaiming: use kaiming punctuations styles for CJK
% - scheme=plain: use plain scheme, do not override \normalsize font size
% - space=auto: space settings for CJK characters
%
% CJK fonts are resolved per locale language, e.g. Noto Serif/Sans CJK JP
% for Japanese resumes, with fallback to the next region if not installed
%
% ref:
% - http://ctan.mirrorcatalogs.com/language/chinese/ctex/ctex.pdf
\usepackage[UTF8, heading=false, punct=kaiming, scheme=plain, space=auto]{ctex}

\IfFontExistsTF{Noto Serif CJK SC}{
  \setCJKmainfont{Noto Serif CJK SC}
}{}
\IfFontExistsTF{Noto Sans CJK SC}{
  \setCJKsansfont{Noto Sans CJK SC}
}{}

%% line spacing
\usepackage{setspace}
\setstretch{1.125}

%% URL styling - use normal text instead of monospace
\urlstyle{same}

%% Patch \httplink and \httpslink to handle full URLs
% The original moderncv macros always prepend http:// or https:// to URLs.
% This causes issues when the URL already has a protocol (e.g., https://example.com)
% resulting in malformed URLs like https://https://example.com.
% This patch checks if the URL already contains "://" and uses it directly if so.
% Note: We use \str_set:Nx (x-expansion) to expand macros like \@homepage first.
\ExplSyntaxOn
\str_new:N \g_yamlresume_url_str

\RenewDocumentCommand{\httpslink}{O{}m}{
  \str_set:Nx \g_yamlresume_url_str {#2}
  \str_if_in:NnTF \g_yamlresume_url_str {://}
    { \url{#2} }
    { \url{https://#2} }
}

\RenewDocumentCommand{\httplink}{O{}m}{
  \str_set:Nx \g_yamlresume_url_str {#2}
  \str_if_in:NnTF \g_yamlresume_url_str {://}
    { \url{#2} }
    { \url{http://#2} }
}
\ExplSyntaxOff

\name{Johannes Fischer}{}
\title{Maschinenbauingenieur – Leichtbau, Simulation und Fahrzeugtechnik}
\phone[mobile]{+49 30 4876 2190}
\email{johannes.fischer@fischer-engineering.de}
\homepage{https://www.johannes-fischer-engineering.de}

\address{Friedrichstraße 112, Berlin, Berlin, Germany, 10117}{}{}

\extrainfo{{\small \faLinkedin}\ \href{https://www.linkedin.com/in/johannes-fischer-engineering}{@johannes-fischer-engineering} {} {} {} • {} {} {} 
{\small \faGithub}\ \href{https://github.com/jfischer-engineering}{@jfischer-engineering}}

\begin{document}

\maketitle

\section{Basics}

\cvline{}{Maschinenbauingenieur mit über sieben Jahren Erfahrung in der Entwicklung, Simulation und Industrialisierung komplexer mechanischer Systeme. Schwerpunkte: Leichtbau, Fahrzeugtechnik, Hydraulik und fertigungsgerechte Konstruktion.
%
\begin{itemize}
\item Sicher im Umgang mit CATIA V5, Siemens NX, ANSYS und MATLAB/Simulink
\item Erfahrung in der Leitung interdisziplinärer Projekte von der Konzeptphase bis zur Serienreife
\item Strukturierte, lösungsorientierte Arbeitsweise mit starkem Qualitäts- und Kostenbewusstsein
\end{itemize}}
\section{Education}

\cventry{Oct 2014–Sep 2017}
        {Master, Maschinenbau – Produktionstechnik, Score: 1,3}
        {Technische Universität Berlin}
        {\url{https://www.tu.berlin}}
        {}
        {\begin{itemize}
\item Vertiefung in Produktionstechnik, Leichtbau und Finite-Elemente-Methoden
\item Masterarbeit zur Topologieoptimierung von Tragstrukturen mit ANSYS
\item Praktische Erfahrung in Laborprojekten zu Werkstoffprüfung und Fertigungsverfahren
\end{itemize}
\textbf{Courses}: Technische Mechanik, Thermodynamik, Strömungsmechanik, Konstruktionslehre, Fertigungsverfahren, Regelungstechnik, Finite-Elemente-Methoden, Werkstoffkunde}

\cventry{Oct 2011–Sep 2014}
        {Bachelor, Maschinenbau, Score: 1,7}
        {Hochschule München}
        {\url{https://www.hm.edu}}
        {}
        {\begin{itemize}
\item Grundlagen in Technischer Mechanik, Thermodynamik und Konstruktionslehre
\item Bachelorarbeit zur Auslegung einer Getriebekomponente
\item Mitglied der Hochschulgruppe Fahrzeugtechnik
\end{itemize}
\textbf{Courses}: Technische Mechanik, Thermodynamik, Konstruktionslehre, Fertigungsverfahren, Werkstoffkunde, Elektrotechnik}
\section{Work}

\cventry{Apr 2021–Present}
        {Senior Entwicklungsingenieur – Fahrzeugtechnik}
        {Nordbahn Mobility GmbH}
        {\url{https://www.nordbahn-mobility.de}}
        {}
        {\begin{itemize}
\item Technische Leitung von Entwicklungsprojekten für Fahrwerks- und Karosseriestrukturen im Schienenfahrzeugbau
\item Konstruktion und Auslegung leichter Tragstrukturen unter Berücksichtigung von Normen (EN 12663, EN 15085)
\item Durchführung von FEM-Berechnungen und Validierung der Ergebnisse durch Tests und Messkampagnen
\item Steuerung externer Zulieferer und Abstimmung mit Fertigung, Einkauf und Qualitätssicherung
\item Mentoring von Nachwuchsingenieuren und Aufbau eines internen CAD-Standards
\end{itemize}
\textbf{Keywords}: Leichtbau, Fahrzeugtechnik, CATIA V5, FEM, EN 12663, Serienentwicklung}

\cventry{Oct 2017–Mar 2021}
        {Entwicklungsingenieur – Hydrauliksysteme}
        {RheinHydraulik Systems AG}
        {\url{https://www.rheinhydraulik.de}}
        {}
        {\begin{itemize}
\item Entwicklung und Optimierung hydraulischer Komponenten und Systeme für mobile Arbeitsmaschinen
\item Erstellung von 3D-Modellen, Fertigungszeichnungen und Stücklisten mit SolidWorks und PDM
\item Durchführung von Toleranzanalysen, Festigkeitsberechnungen und Design-FMEA
\item Begleitung von Prototypenbau, Prüfstandsversuchen und Serienanlauf
\item Identifikation von Kosteneinsparungen durch konstruktive Optimierungen und Werkstoffauswahl
\end{itemize}
\textbf{Keywords}: Hydraulik, SolidWorks, PDM, Design-FMEA, Serienanlauf, Kostenoptimierung}
\section{Languages}

\cvline{German}{Native or Bilingual Proficiency \hfill \textbf{Keywords}: Muttersprache}
\cvline{English}{Full Professional Proficiency \hfill \textbf{Keywords}: TOEFL iBT 108, Technisches Englisch}
\cvline{Spanish}{Limited Working Proficiency \hfill \textbf{Keywords}: Grundkenntnisse, Reiseerfahrung}
\section{Skills}

\cvline{CAD-Konstruktion}{Expert \hfill \textbf{Keywords}: CATIA V5, Siemens NX, SolidWorks, AutoCAD, PDM/PLM}
\cvline{Berechnung \& Simulation}{Advanced \hfill \textbf{Keywords}: FEM, ANSYS, Abaqus, CFD, MATLAB/Simulink}
\cvline{Fertigung \& Produktion}{Advanced \hfill \textbf{Keywords}: Drehen, Fräsen, Schweißen, Lean Production, Qualitätssicherung}
\cvline{Projektmanagement}{Intermediate \hfill \textbf{Keywords}: Agile Methoden, Meilensteinplanung, Kostenkontrolle, Stakeholder}
\cvline{Programmierung \& Datenanalyse}{Intermediate \hfill \textbf{Keywords}: Python, VBA, SQL, Datenvisualisierung}
\section{Awards}

\cventry{Mar 2018}
        {VDI – Verein Deutscher Ingenieure}
        {VDI Förderpreis für hervorragende Abschlussarbeiten}
        {}
        {}
        {Auszeichnung für die Masterarbeit zur Topologieoptimierung von Leichtbaustrukturen im Schienenfahrzeugbau.}

\cventry{Nov 2017}
        {Technische Universität Berlin}
        {Preis für die beste Masterarbeit}
        {}
        {}
        {Auszeichnung des Fachgebiets Konstruktion und Produktionstechnik für herausragende wissenschaftliche Arbeit.}
\section{Certificates}

\cventry{Jun 2019}
        {Dassault Systèmes}
        {Certified SolidWorks Professional (CSWP)}
        {\url{https://www.solidworks.com/certification}}
        {}
        {}

\cventry{Sep 2020}
        {TÜV Rheinland}
        {Six Sigma Green Belt}
        {\url{https://www.tuv.com}}
        {}
        {}

\cventry{May 2022}
        {IPMA}
        {IPMA Level D – Projektmanagement}
        {\url{https://www.ipma.world}}
        {}
        {}
\section{Publications}

\cventry{Oct 2020}
        {Leichtbauoptimierung von Schienenfahrzeugstrukturen mittels parametrischer FEM}
        {VDI-Verlag}
        {\url{https://www.vdi.de/verlag}}
        {}
        {\begin{itemize}
\item Untersuchung parametrischer FEM-Modelle zur Gewichtsreduktion von Schienenfahrzeugstrukturen
\item Vergleich von Aluminium- und Faserverbundwerkstoffen hinsichtlich Steifigkeit, Kosten und Fertigbarkeit
\item Handlungsempfehlungen für die frühe Entwicklungsphase im Schienenfahrzeugbau
\end{itemize}}
\section{Projects}

\cventry{Jan 2022–Nov 2023}
        {Entwicklung eines hybriden Tragrahmens aus Aluminium und Faserverbundwerkstoffen zur Reduktion des Fahrzeuggewichts bei gleichbleibender Steifigkeit.}
        {Hybrid-Leichtbau-Tragrahmen für Schienenfahrzeuge}
        {\url{https://www.nordbahn-mobility.de/projekte/leichtbau-tragrahmen}}
        {}
        {\begin{itemize}
\item Konzeptentwicklung und Variantenbewertung in enger Abstimmung mit Fahrzeugarchitektur und Fertigung
\item Parametrische FEM-Modelle zur Auslegung und Optimierung der Wandstärken und Fügepunkte
\item Planung und Begleitung von Prüfstandsversuchen zur Validierung der Auslegung
\item Dokumentation der Ergebnisse und Vorbereitung der Serienentscheidung
\end{itemize}
\textbf{Keywords}: Hybrid-Leichtbau, Faserverbund, FEM}
\end{document}